\section{Practical Test Expectations}

This section explains how the practical-test-related tutorial submissions are used in the module. You must submit all assigned weekly tutorial exercises through the relevant ITEL tabs, even though only selected activities are graded formally. These submissions help the lecturer verify your individual contribution and your understanding of the methods introduced in class.

The tutorial sessions that receive particular assessment attention are:
\begin{itemize}
    \item \textbf{Identifying Customer Needs};
    \item \textbf{Product Specifications}; and
    \item \textbf{Concept Generation}.
\end{itemize}

Submissions must be uploaded via the respective weekly tabs on ITEL.

Although only three tutorials will be graded, the other submissions are still important. They will be used to verify your individual contributions and may be referenced during the presentation session. Therefore, ensure that all tutorial work is completed diligently and submitted on time.





\subsection{What the Lecturer Will Look For}

The criteria below show how the submitted work supports the relevant \textbf{Course Learning Outcomes (CLOs)}.

\subsubsection{Identifying Customer Needs}
\begin{itemize}
    \item Clarity of Survey Objectives
    \item Quality of Survey Questions and Prompts
    \item Accuracy of Needs Statement Translation
    \item Quantity of Translated Needs per Theme
    \item Number of Translated Themes
    \item Classification of Need Types
    \item Source and Relevance of Customer Statements
\end{itemize}

\subsubsection{Product Specifications}
\begin{itemize}
    \item Importance vs. Needs Table
    \item Functionality-Focused Metrics
    \item Inclusion of Additional Metrics for Key Needs
    \item Justification for Metric Units
    \item Needs-to-Metric Matrix
    \item Competitive Benchmarking Chart (CBC) Metrics
    \item Sources Used for CBC Comparison
    \item Alignment Between CBC and Identified Needs
    \item Bar Plot Comparing Group Concept with Competitors
    \item Proper Table and Figure Captions
\end{itemize}

\subsubsection{Concept Generation}
\begin{itemize}
    \item Analysis of Primary Functions
    \item Justification for Problem Decomposition Methods
    \item Use of Action Verbs in Function Representation
    \item Quality of Proposed Solutions
    \item Justification for Selected Solutions
    \item Quantity of Proposed Solutions
    \item Variety of Information Sources
    \item Quality of Citations
\end{itemize}

% \subsection{Reflection Prompt}

% After each practical-test-related submission, briefly ask yourself: Which part of the method do I understand well, and which part still needs stronger evidence, better structure, or clearer explanation?

\subsection{Latex or Word?}
In last semester, I enforce the use of a LaTeX template for all tutorial submissions. The reason is that LaTeX encourages clear structure, proper formatting, and consistent presentation of technical content. It also allows for better integration of figures, tables, and references, which are essential for communicating engineering ideas effectively. Apart from that, using the combination of Latex and LLMs can help you to focus on the content of your work. This is because, when copy-pasting from LLM, you no need to worry about formatting, and you can just focus on the content. The Laxex template can be accessed through the following link: \href{https://itel.ums.edu.my/mod/resource/view.php?id=968033}{Download Me}.

However, I understand that some students may not be familiar with LaTeX. Therefore, for this semester, you are allowed to use any word processing software you are comfortable with, as long as your submission follows the expected structure.
